% GENERATED FILE. Edit canonical lessons, then run npm run generate:books.
% canonical-source-hash: fnv1a64:e2f64b84c98ed3c1
% canonical-lessons: SA-C06-numbers-1-5, SA-C06-numbers-1-5-more, SA-C06-number-cognates, SA-S206-letter-nya, SA-C06-number-cognates-more, SA-C06-pancha-travels, SA-C06-pancha-travels-more, SA-R06-second-pass-where-the-words-came-from

\chapter{Numbers One to Five}
\label{ch:numbers-one-to-five}
\hlchaptermodality{6}

\begin{chapteropening}
\textbf{By the end of this chapter:} \emph{I can say the Sanskrit numerals \sk{एक} to \sk{पञ्च} with their dual and gendered forms, and follow \sk{पञ्च} outward into Punjab, pentagon, and the disputed history of punch.}

Trace pañca out of Sanskrit: into Persian panj-āb, the five waters behind Punjab; into Greek pente in pentagon and pentathlon; and into punch, with the five-ingredients account labelled as common rather than settled. Then count one to five with the dual and the neuter forms the daughter languages inherit, say the cousins and the sound laws between them, and pick out \sk{ञ}.
\end{chapteropening}

\section[\texorpdfstring{\sk{एक} \sk{द्व} \sk{त्रि} \sk{चतुर्} \sk{पञ्च} (eka dva tri …)}{एक द्व त्रि चतुर् पञ्च (eka dva tri …)}]{\sk{एक}, \sk{द्व}, \sk{त्रि}, \sk{चतुर्}, \sk{पञ्च} — one to five}

\label{lesson:SA-C06-numbers-1-5}



\begin{quote}
This chapter does something the other Indic chapters can't do alone. Sanskrit preserves the \textbf{Old Indo-Aryan forms behind} the Hindi, Marathi, Bengali, Gujarati and Punjabi ones — and, at the same time, these forms are \textbf{cousins} of the English ones. Learn five and you get a foothold in both directions.

How many numbers does a Sanskrit verb count in? (Three: singular, dual for exactly two, and plural.)
\end{quote}

\subsection*{You'll want to know: the five}
\noindent
\begin{tabularx}{\linewidth}{@{}>{\raggedright\arraybackslash}X>{\raggedright\arraybackslash}X>{\raggedright\arraybackslash}X>{\raggedright\arraybackslash}X@{}}
\toprule
\textbf{} & \textbf{stem} & \textbf{masculine} & \textbf{\textbf{neuter}} \\
\midrule
1 & \emph{éka-} & \textbf{\sk{एकः}} \emph{ékaḥ} & \sk{एकम्} \emph{ékam} \\
2 & \emph{dvá-} & \emph{dváu} & \emph{dvé} \\
3 & \emph{tri-} & \textbf{\sk{त्रयः}} \emph{tráyaḥ} & \emph{tr\'{\={\i}}ṇi} \\
4 & \emph{catúr-} & \textbf{\sk{चत्वारः}} \emph{catv\'{\={a}}raḥ} & \textbf{\sk{चत्वारि}} \emph{catv\'{\={a}}ri} \\
5 & \sk{पञ्च} \emph{páñca} & — & — \\
\bottomrule
\end{tabularx}

Two notes first. A third, about the last column, follows in the next lesson.

\textbf{Sanskrit has a dual} — a separate number for exactly two things — which is why \textquotedblleft{}two\textquotedblright{} is \emph{dváu} rather than a plural.

\textbf{\emph{Pañca} barely inflects} in the nominative and accusative. It is also the one you already half-know.

\begin{grammarlens}[title={What you've built: two directions}]
These five forms point in two directions. The next lesson lines them up with Latin and English and explains the sound laws. A final short lesson follows \emph{pañca} into \emph{Punjab}, \emph{pentagon}, and the qualified history of \emph{punch}.
\end{grammarlens}

\subsection*{Your turn}
Say these aloud:

\begin{itemize}
\raggedright
  \item \textquotedblleft{}eka, dva, tri, catur, pañca\textquotedblright{}
  \item the exactly-two category — dual
\end{itemize}

\begin{tcolorbox}[breakable,colback=teal!4,colframe=teal!35!black,title={Before you move on}]
Give one to five in Sanskrit. (\emph{Eka, dva, tri, catur, pañca}.) Why does \textquotedblleft{}two\textquotedblright{} look unusual? (Sanskrit has a \textbf{dual} — a form for exactly two.) Which number barely inflects in this table? (\emph{Pañca}.) Next: connect the Sanskrit forms west to Latin and English.
\end{tcolorbox}
\section[\texorpdfstring{(\sk{एक} \sk{द्व} \sk{त्रि} \sk{चतुर्} \sk{पञ्च}, continued)}{(एक द्व त्रि चतुर् पञ्च, continued)}]{(\sk{एक} \sk{द्व} \sk{त्रि} \sk{चतुर्} \sk{पञ्च}, continued) — numbers have gender}

\label{lesson:SA-C06-numbers-1-5-more}



\begin{quote}
Before going on: count one to five again — \textbf{\sk{एक} \sk{द्व} \sk{त्रि} \sk{चतुर्} \sk{पञ्च}}.

And the letter from the last lesson: say the sound of \sk{अ}.

From the goodbyes: which English word is a cousin of \emph{gam} \textquotedblleft{}to go\textquotedblright{}? (\emph{Come}.)
\end{quote}

\subsection*{You'll want to know: numbers have gender}
\textbf{Sanskrit numbers have gender.} Look back at the last column of the table in the last lesson. These same words, in Hindi, Marathi, Bengali, Gujarati, and Punjabi, mostly descend from the \textbf{neuter} forms — \emph{tr\'{\={\i}}ṇi} and \emph{catv\'{\={a}}ri}, not \emph{tráyaḥ} and \emph{catv\'{\={a}}raḥ}. Keep that column in view; the daughter languages inherit from it.

\subsection*{Your turn}
\begin{itemize}
\raggedright
  \item \emph{Say it:} what this lesson added to \textbf{\sk{एक} \sk{द्व} \sk{त्रि} \sk{चतुर्} \sk{पञ्च}}, in one sentence
\end{itemize}

\begin{tcolorbox}[breakable,colback=teal!4,colframe=teal!35!black,title={Before you move on}]
What did this lesson add to \textbf{\sk{एक} \sk{द्व} \sk{त्रि} \sk{चतुर्} \sk{पञ्च}}? Say it once more, without looking.
\end{tcolorbox}
\section[\texorpdfstring{\sk{एक} · \sk{द्व} · \sk{त्रि} … (eka · dva · tri …)}{एक · द्व · त्रि … (eka · dva · tri …)}]{The Sanskrit numbers are English cousins}

\label{lesson:SA-C06-number-cognates}



\begin{quote}
Give the five stems: \emph{eka, dva, tri, catur, pañca}.
\end{quote}

\begin{grammarlens}[title={What you've built: one family, not borrowing}]
\noindent
\begin{tabularx}{\linewidth}{@{}>{\raggedright\arraybackslash}X>{\raggedright\arraybackslash}X>{\raggedright\arraybackslash}X@{}}
\toprule
\textbf{Sanskrit} & \textbf{Latin} & \textbf{English} \\
\midrule
\emph{éka-} & \emph{ūnus} & \textbf{one} \\
\emph{dvá-} & \emph{duo} & \textbf{two} \\
\emph{tri-} & \emph{trēs} & \textbf{three} \\
\emph{catúr-} & \emph{quattuor} & \textbf{four} \\
\emph{páñca} & \emph{quīnque} & \textbf{five} \\
\bottomrule
\end{tabularx}

Sanskrit, Latin, and English inherited these from one family. Numbers resist replacement, so their relationship remains audible after millennia.

The first row needs one caveat: \emph{éka-} and \emph{ūnus/one} use the same root with different suffixes, PIE *\emph{oy-ko-} versus *\emph{óynos}. They are relatives rather than the same complete word. The other four rows continue the same word.
\end{grammarlens}

\subsection*{Your turn}
Say these aloud:

\begin{itemize}
\raggedright
  \item \emph{tri, trēs, three}
  \item regular inheritance or borrowing — inheritance
\end{itemize}

\begin{tcolorbox}[breakable,colback=teal!4,colframe=teal!35!black,title={Before you move on}]
Are the Latin and Sanskrit forms borrowed from one another? (No: they are inherited cousins.) Next: the sound laws that pulled the cousins apart.
\end{tcolorbox}
\section[\texorpdfstring{\sk{ञ} (ña)}{ञ (ña)}]{\sk{ञ} — one character, met inside words you already say}

\label{lesson:SA-S206-letter-nya}



\begin{quote}
Before the new one, the ones you have already met: \sk{य} · \sk{◌ृ} · \sk{अ}. Say what each of them does.

One character this time. One only — and you have been saying it for pages without knowing which mark on the page it was.
\end{quote}

\subsection*{Script you'll notice: \sk{ञ}}
\textbf{\sk{ञ}} — \emph{ña}.

It is a \textbf{consonant}, and in this script a consonant is never bare: it comes with an \emph{a} already in it. So it is not \emph{ñ}, it is \textbf{ña}.

What it is made of:

\begin{itemize}
\raggedright
  \item a clockwise open-left bowl
  \item a rightward shoulder rising to the headline
  \item a short descending lower stem
  \item the top shirorekhā
\end{itemize}

You already say these words, and \sk{ञ} is one of the shapes inside them — the rest of their shapes are still ahead of you, so here the character stands on its own:

\begin{itemize}
\raggedright
  \item \emph{eka dva tri catur pañca} — one to five — Old Indo-Aryan ancestors preserved in Sanskrit, and cousins of the English ones
  \item \emph{eka · dva · tri · catur · pañca} — the Sanskrit, Latin, and English number family and the sound laws between them
\end{itemize}

\subsection*{Writing: \sk{ञ}}
\hlblockfigure{\detokenize{figures/SA-S206-letter-nya-filmstrip.pdf}}{How \sk{ञ} is written, stroke by stroke}

\begin{itemize}
\raggedright
  \item \textbf{1.} start at the open upper-left tip and curve clockwise around the bowl to the open lower-left tip
  \item \textbf{2.} lift at the bowl's right junction, sweep right through the shoulder, then curve upward along the right side to the headline
  \item \textbf{3.} lift at the lower right-side junction and descend the short stem below the bowl
  \item \textbf{4.} lift and draw the top shirorekhā left-to-right
\end{itemize}

\textbf{Pen lifts: 3.} The pen comes up 3 times and no more.

\begin{quote}
verified four-stroke teaching form fitted to the bundled printed outline
\end{quote}

\begin{quote}
This is one attested teaching order and not a national standard — handwriting here is taught with school-to-school variation. Source: Opiaterein, ‘Deva-\sk{ञ}-order.gif’, strokes 1–4, Wikimedia Commons, 10 May 2009.
\end{quote}

\subsection*{Your turn}
\begin{itemize}
\raggedright
  \item \emph{Look:} at this, and find \sk{ञ} in it
\end{itemize}

\begin{quote}
\sk{ञ}
\end{quote}

\begin{itemize}
\raggedright
  \item \emph{Trace it:} \sk{ञ} three times, saying \emph{ña} as you finish each one
  \item \emph{Look:} back at any page of this chapter and find \sk{ञ} once more
\end{itemize}

\begin{tcolorbox}[breakable,colback=teal!4,colframe=teal!35!black,title={Before you move on}]
Which character is this — \sk{ञ}? What sound does it carry? (\emph{\textbf{ña}}.) Name one word you already say that contains it.
\end{tcolorbox}
\section[\texorpdfstring{(\sk{एक} · \sk{द्व} · \sk{त्रि} · \sk{चतुर्} · \sk{पञ्च}, continued)}{(एक · द्व · त्रि · चतुर् · पञ्च, continued)}]{(\sk{एक} · \sk{द्व} · \sk{त्रि} · \sk{चतुर्} · \sk{पञ्च}, continued) — The word, taken apart — rounded \emph{k} goes three ways}

\label{lesson:SA-C06-number-cognates-more}



\begin{quote}
Before going on: say \textbf{\sk{एक} · \sk{द्व} · \sk{त्रि} · \sk{चतुर्} · \sk{पञ्च}} again, and what it means.
\end{quote}

\begin{cousinweb}
PIE *\emph{k\textsuperscript{w}} was a \emph{k} with rounded lips. In *\emph{k\textsuperscript{w}etwóres} “four,” three branches treated it differently:

\noindent
\begin{tabularx}{\linewidth}{@{}>{\raggedright\arraybackslash}X>{\raggedright\arraybackslash}X>{\raggedright\arraybackslash}X@{}}
\toprule
\textbf{branch} & \textbf{outcome} & \textbf{example} \\
\midrule
Latin & kept \emph{qu-} & \emph{\textbf{qu}attuor} \\
Indo-Iranian & merged it with \emph{k}, then palatalised before front vowels & \emph{\textbf{c}atvāri} \\
Germanic & shifted it to \emph{hw-} & \emph{\textbf{wh}at}, \emph{\textbf{wh}o} \\
\bottomrule
\end{tabularx}

Sanskrit \emph{c-} is therefore a regular result, not a random replacement. The same sequence helps explain \emph{pañca}'s \emph{ñc}.
\end{cousinweb}

\begin{cousinweb}
The \textbf{f-} of \emph{five} does not come from *\emph{k\textsuperscript{w}}. It comes from initial *\emph{p}, changed to \emph{f} by \textbf{Grimm's law}, as in \emph{pater $\to$ father} and \emph{pēs $\to$ foot}.

English \emph{four} is irregular. By the sound rule it should begin \emph{hw-} like \emph{what}. Its \emph{f-} is usually explained as analogy with neighbouring \emph{five}. Counting words can reshape one another.
\end{cousinweb}

\subsection*{Your turn}
Say these aloud:

\begin{itemize}
\raggedright
  \item rounded \emph{k} — \emph{quattuor, catvāri, what}
  \item initial \emph{p $\to$ f} — \emph{pañca/five}
\end{itemize}

\begin{tcolorbox}[breakable,colback=teal!4,colframe=teal!35!black,title={Before you move on}]
What are the three outcomes of *\emph{k\textsuperscript{w}}? (Latin \emph{qu-}, Indo-Iranian \emph{c-} after merger and palatalisation, Germanic \emph{hw-}.) Where does \emph{five}'s initial \emph{f} come from? (PIE \emph{p} by Grimm's law.) What is the usual explanation for \emph{four} also beginning with \emph{f}? (Analogy with neighbouring \emph{five}.)
\end{tcolorbox}
\section[\texorpdfstring{\sk{पञ्च} (pañca)}{पञ्च (pañca)}]{\emph{Pañca} hiding in familiar words}

\label{lesson:SA-C06-pancha-travels}



\begin{quote}
Say Sanskrit “five.” (\emph{Pañca}.) Give its English cousin. (\emph{Five}.)
\end{quote}

\begin{cousinweb}
\textbf{Punjab} comes from Persian \emph{panj-āb}, “five waters,” for the land of five rivers spanning India and Pakistan. Persian \emph{panj} is the Iranian cousin of Indic \emph{pañca}: both descend from the same Indo-Iranian and PIE number root.

The place-name is Persian, not Sanskrit, but the family resemblance is real.
\end{cousinweb}

\begin{cousinweb}
Greek \emph{\textbf{pente}} is another descendant of the same root. It appears in \textbf{pentagon}, a five-sided figure, and \textbf{pentathlon}, a five-event contest.
\end{cousinweb}

\subsection*{Your turn}
Say these aloud:

\begin{itemize}
\raggedright
  \item \emph{panj + āb} — five waters
  \item Greek \emph{pente} — pentagon, pentathlon
\end{itemize}

\begin{tcolorbox}[breakable,colback=teal!4,colframe=teal!35!black,title={Before you move on}]
What does \emph{Punjab} mean? (“Five waters.”) Is Persian \emph{panj} borrowed from Sanskrit? (No: it is an Iranian cousin.) Which Greek form hides in \emph{pentagon}? (\emph{Pente}.) Next: the drink \emph{punch}, and why its story comes with a caution.
\end{tcolorbox}
\section[\texorpdfstring{(\sk{पञ्च}, continued)}{(पञ्च, continued)}]{(\sk{पञ्च}, continued) — The word, taken apart — the qualified \emph{punch} story}

\label{lesson:SA-C06-pancha-travels-more}



\begin{quote}
Before going on: say \textbf{\sk{पञ्च}} again, and what it means.
\end{quote}

\begin{cousinweb}
The usual account derives \textbf{punch}, the drink, from Hindi \emph{pāṁch} $\leftarrow$ Sanskrit \emph{pañca}, because the drink was said to contain five ingredients.

That origin is not unanimous. A rival derives the word from \textbf{puncheon}, a type of cask. Keep the useful connection with a label: “the common five-ingredients account,” not an unquestioned fact.
\end{cousinweb}

\subsection*{Your turn}
Say these aloud:

\begin{itemize}
\raggedright
  \item \emph{punch} — common account, not certain
  \item the rival: \emph{puncheon}, a cask
\end{itemize}

\subsection*{Your turn — the chapter's five, and what they carry}
Say these aloud:

\begin{itemize}
\raggedright
  \item one to five (\emph{eka, dva, tri, catur, pañca})
  \item why \textquotedblleft{}two\textquotedblright{} is \emph{dváu}, not a plural (a \textbf{dual}, for exactly two)
  \item the two forms of \textquotedblleft{}three\textquotedblright{}, and which one Hindi, Marathi and the rest mostly inherit (\emph{tráyaḥ}, masculine; \emph{tr\'{\={\i}}ṇi}, the \textbf{neuter})
  \item the two directions these five point (back to the forms behind the modern Indic numbers; across to the English ones, as cousins)
\end{itemize}

\subsection*{Your turn — cousins and sound laws}
Say these aloud:

\begin{itemize}
\raggedright
  \item \textquotedblleft{}three\textquotedblright{} in Sanskrit, Latin and English, and whether one borrowed it (\emph{tri, trēs, three}; no — inherited cousins)
  \item why \emph{éka} and \textbf{one} are relatives but not one whole word (same root, different suffixes)
  \item what rounded \emph{*k\textsuperscript{w}} became in Latin, Sanskrit and English (\emph{qu-}, \emph{c-}, \emph{hw-}: \emph{quattuor}, \emph{catvāri}, \textbf{what})
  \item where the \emph{f-} of \textbf{five} comes from, and why \textbf{four} has one too (PIE \emph{*p}, by Grimm's law; analogy with \emph{five})
\end{itemize}

\subsection*{Your turn — Punjab and pentagon}
Say these aloud:

\begin{itemize}
\raggedright
  \item what \textbf{Punjab} means, and whether Persian \emph{panj} is a loan from Sanskrit (\emph{panj-āb}, \textquotedblleft{}five waters\textquotedblright{}; no — an Iranian cousin)
  \item the Greek five, and two English words built on it (\emph{pente}; \textbf{pentagon}, \textbf{pentathlon})
\end{itemize}

\subsection*{Script — the shape inside \emph{pañca}}
\begin{quote}
\sk{अ} · \sk{ञ}
\end{quote}

\begin{itemize}
\raggedright
  \item \emph{Point to:} \emph{ña}, then \emph{a}
  \item \emph{Say it:} the sound \sk{ञ} carries, and the number it sits inside (\emph{ña}; \emph{pañca})
\end{itemize}

\begin{tcolorbox}[breakable,colback=teal!4,colframe=teal!35!black,title={Before you move on}]
What caution belongs with the drink \emph{punch}? (The five-ingredients derivation is common, but \emph{puncheon} is a rival account.)
\end{tcolorbox}
\section[\texorpdfstring{\sk{नमस्कारः} (namaskāraḥ)}{नमस्कारः (namaskāraḥ)}]{\sk{नमस्कारः} — the ending you can hear, and four claims you were told once}

\label{lesson:SA-R06-second-pass-where-the-words-came-from}



\begin{quote}
Say the longer of the two greetings. Then say what its last breath — the soft \emph{-ḥ} — is telling you about the word.

Sanskrit numbers have gender. Which forms of \textquotedblleft{}three\textquotedblright{} and \textquotedblleft{}four\textquotedblright{} do the modern languages mostly inherit? (The neuter ones, \emph{trīṇi} and \emph{catvāri}.)
\end{quote}

\begin{culture}
\emph{Namaskāraḥ} ends on a \textbf{visarga}, and a visarga on a noun like this one is not decoration: it marks the \textbf{masculine singular}. You have been hearing it since your first pair of greetings without being asked to notice it twice.

The numbers in front of you now make the point differently. \emph{Pañca} has no visarga; \emph{namaskāraḥ} does. One is a bare stem, the other is a word doing a job in a sentence. Sanskrit tells you which by the ending, not by the position.
\end{culture}

\begin{tcolorbox}[breakable,skin=enhanced,colback=violet!6,colframe=violet!45!black,boxrule=0.5pt,arc=1mm,left=6pt,right=6pt,top=4pt,bottom=4pt,fonttitle=\bfseries,title={Across the family}]
Four historical claims have gone past you, one at a time, each at the end of a chapter. Set against each other they stop being trivia and become one argument:

\noindent
\begin{tabularx}{\linewidth}{@{}>{\raggedright\arraybackslash}X>{\raggedright\arraybackslash}X@{}}
\toprule
\textbf{the claim} & \textbf{what it rests on} \\
\midrule
a judge in Calcutta announced in 1786 that Sanskrit, Greek and Latin came from a common source & he was comparing verbs and grammar, not lists of similar-sounding words \\
the words of the naming exchange are traceable to that same source & \emph{nāma}, \emph{mama}, \emph{tvam}, \emph{kim} all have cousins west of India \\
the wellbeing words have roots reaching east and west alike & the root is the unit that travels, not the finished word \\
the leave-taking words are built from roots you can already name & \emph{punar-darśanāya} is \emph{again} plus \emph{for seeing} \\
\bottomrule
\end{tabularx}

Read the right-hand column downward. Every one of the four is an argument from \textbf{roots and endings}, which is the same evidence the numbers in this chapter rest on. That is why the claim about Sanskrit and Latin is a claim about grammar and not about vocabulary.
\end{tcolorbox}

\begin{tcolorbox}[breakable,colback=teal!4,colframe=teal!35!black,title={Before you move on}]
What kind of evidence did the 1786 announcement rest on — matching words, or matching grammar? (\textbf{Grammar.}) What does the visarga on \emph{namaskāraḥ} mark? (\textbf{Masculine singular.})
\end{tcolorbox}
